\section{Preliminarii}

\subsection{Funcția hash}

O funcție hash criptografică este o transformare deterministă $h \colon \{0,1\}^* \to \{0,1\}^n$ care asociază datelor de intrare de lungime arbitrară un hash de lungime fixă $n$.
Pentru a fi utilă în context criptografic, ea trebuie să fie eficientă din punct de veder computațional și să satisfacă trei proprietăți de securitate ~\cite{merkle1989, damgard1989}: 
\begin{enumerate}
    \item Rezistența la preimagine: dat fiind un hash $y$, este infezabil să se găsească un mesaj $m$ cu $h(m) = y$.
    \item Rezistența la a doua preimagine: dat un mesaj $m$, este infezabil să se găsească $m' \neq m$ cu $h(m') = h(m)$.
    \item Rezistența la coliziuni: este infezabil să se găsească orice pereche $m \neq m'$ cu $h(m) = h(m')$.
\end{enumerate}

\subsection{Sistemul dinamice}

Un sistem dinamic discret este un triplet $(T, X, f)$, unde:
\begin{enumerate}
    \item $T$ este spațiul timpului, iar $(T, +)$ este un monoid aditiv
    \item $X$ este spațiul de stări sau spațiul fazelor
    \item $f \colon T \times X \to X$ este funcția de evoluție, care satisface următoarele condiții:
    \begin{enumerate}
        \item $f(0, x) = x$
        \item $f(t+s, x) = f(t, f(s, x))$ pentru orice $t, s \in T$ și $x \in X$
    \end{enumerate}
\end{enumerate}
Devaney ~\cite{devaney1989} definește un sistem haotic ca având trei proprietăți esențiale: sensibilitate la condițiile inițiale, tranzitivitate topologică și faptul că mulțimea orbitelor periodice este densă. 
Banks ~\cite{banks1992} demonstrează că tranzitivitatea topologică și densitatea orbitelor periodice implică sensibilitatea la condițiile inițiale. \\

Exponentul Lyapunov măsoară rata medie de separare exponențială a două orbite pornite din condiții inițiale infinitezimal de apropiate.
Pentru un sistem unidimensional cu funcția de evoluție $f$, el este definit ca:
\begin{equation}
    \lambda = \lim_{N \to \infty} \frac{1}{N} \sum_{k=0}^{N-1} \ln \abs{f'(x_k)},
\end{equation}
unde $x_k$ sunt iteratele orbitei.
O valoare $\lambda > 0$ indică divergența exponențială a orbitelor vecine și constituie semnătura cantitativă a comportamentului haotic, în timp ce $\lambda \le 0$ corespunde unei dinamici stabile sau cvasi-periodice~\cite{devaney1989, alvarez2006}. \\

Relevanța exponentului Lyapunov pentru o funcție hash reiese tocmai din această sensibilitate la condițiile inițiale.
Un exponent strict pozitiv garantează că modificarea unui singur bit din datele de intrare se propagă exponențial pe parcursul iterațiilor, separând rapid orbitele corespunzătoare a două mesaje aproape identice.
Acest mecanism stă la baza efectului de avalanșă dorit, prin care date de intrare cu diferențe minime produc hash-uri complet necorelate, iar magnitudinea exponentului oferă o măsură cantitativă a vitezei cu care difuzia atinge saturația~\cite{alvarez2006}. \\

\subsection{Construcția Sponge}

Construcția Sponge este o schemă de proiectare a funcțiilor hash care operează asupra unei stări interne de lățime $b$ biți.
Starea internă este împărțită într-o rată $r$ și o capacitate $c$.
Asupra stării interne se aplică în mod repetat o permutare fixă $f$~\cite{bertoni2007sponge}.
Procesarea se desfășoară în două faze: în faza de absorbție, blocurile mesajului, de câte $r$ biți, sunt XOR-ate cu partea de rată a stării.
Ulterior se aplică permutarea $f$. 
În faza de stoarcere, biții pentru hash sunt extrași din partea de rată, intercalați cu aplicări succesive ale lui $f$, până la obținerea lungimii dorite a hash-ului.
Capacitatea $c$, care nu este expusă niciodată direct la intrare sau la ieșire, determină nivelul de securitate al construcției. \\

\begin{figure}[ht]
    \centering
    \includegraphics[width=0.8\textwidth]{sponge.png}
    \caption{Construcția Sponge ~\cite{bertoni2007sponge}.}
    \label{fig:sponge}
\end{figure}

\subsection{Arborele Merkle}

Arborele Merkle este o structură de date arborescentă, de regulă binară, în care fiecare frunză conține hash-ul unui bloc de date.
Fiecare nod intern conține hash-ul obținut prin concatenarea și aplicarea funcției hash asupra hash-urilor copiilor săi~\cite{merkle1989}.
Hash-ul nodului rădăcină, numit rădăcină Merkle, sintetizează astfel întregul set de date într-o singură valoare de lungime fixă.
Orice modificare a unui bloc se propagă în mod inevitabil până la rădăcină.
Această organizare permite paralelizarea funcțiilor hash, întrucât subarborii independenți pot fi procesați simultan. \\
